\subsection{Introducción}

En la actualidad, el uso de tecnologías digitales se ha convertido en un elemento fundamental para el desarrollo y crecimiento de los negocios. La transformación digital ha permitido que empresas y emprendimientos adopten nuevas herramientas que facilitan la comunicación con los clientes, optimizan los procesos de atención y mejoran la eficiencia en la gestión de la información. En este contexto, las aplicaciones de mensajería instantánea como WhatsApp y Telegram se han consolidado como uno de los principales canales de interacción entre empresas y usuarios, debido a su facilidad de uso, accesibilidad y amplia presencia en la vida cotidiana de las personas.

Sin embargo, muchos negocios aún gestionan la atención a sus clientes de forma manual, lo que puede generar retrasos en las respuestas, saturación de mensajes y dificultades para brindar información oportuna. Esta situación se vuelve más evidente cuando el volumen de consultas aumenta, provocando una disminución en la calidad del servicio y, en algunos casos, la pérdida de posibles clientes.

Ante esta problemática, el desarrollo de sistemas automatizados basados en bots conversacionales representa una alternativa eficiente para mejorar la atención digital. Estas herramientas permiten automatizar respuestas frecuentes, brindar información inmediata y mantener una comunicación constante con los usuarios, reduciendo la carga operativa y optimizando el tiempo de respuesta.

El presente proyecto propone el desarrollo un bot automatizado integrado a las plataformas de WhatsApp y Telegram, diseñado para responder consultas de manera inmediata y facilitar la interacción entre los negocios y sus clientes. El sistema permitirá gestionar mensajes automáticos, proporcionar información relevante y apoyar en la atención digital, contribuyendo a mejorar la experiencia del usuario y la eficiencia en la comunicación empresarial.